%% Qualifying-paper template, UHM Linguistics.
%% Compile with LuaLaTeX (or XeLaTeX) + Biber. latexmk does it all:
%%     latexmk main.tex
%% On Overleaf the included latexmkrc selects LuaLaTeX automatically.
%%
%% Change the journal style with profile=  (langsci, ldc, ol, stuf, blank)
%% Change the font with font=              (charis, libertinus, gentium, stix, noto)
%% Omit font= to use the profile's own font.
\documentclass[
  profile = langsci,
  % font  = charis,
  bibfile = references.bib,
  % draft,          % date stamp and overfull-line marks
  % linenumbers,    % line numbers for committee comments
]{uhmqp}

%% Your own gloss abbreviations (beyond the Leipzig set), e.g.
%% \newleipzig{evid}{evid}{evidential}   ->  \Evid
\newleipzig{rdp}{rdp}{reduplication}

\title{Title of the qualifying paper in sentence case}
\shorttitle{Short title}
\author{Author Name}
\affiliation{University of Hawaiʻi at Mānoa}
\email{author@hawaii.edu}
\abstract{Abstract text (100--300 words). Placeholder.}
\keywords{keyword one; keyword two; keyword three}
\qpnote{Qualifying paper 1. Department of Linguistics, University of Hawaiʻi at Mānoa. Submitted Month Year.}

\begin{document}
\maketitle

\section{Introduction}\label{sec:intro}
Placeholder text with a citation \parencite[514--515]{murray-vennemann1983} and a
narrative citation \textcite[238]{mccarthy-prince1999}. An object-language form
\form{ta-wa} \gloss{go} and a grapheme \grapheme{q}.\footnote{A footnote.}
See \cref{sec:data}, \cref{ex:basic}, \cref{tab:sample} and \cref{fig:sample}.

\section{Data}\label{sec:data}
\subsection{Glossed examples}

\ea\label{ex:basic}
  \gll ta-wa ku\\
       \Fsg-go house\\
  \glt `I go home.' \exsource{Source 2020: 12}
\z

\ea\label{ex:sub}
  \ea\label{ex:sub-a}
    \gll ta-wa ku\\
         \Fsg-go house\\
    \glt `I go home.'
  \ex\label{ex:sub-b}
    \gll ta\textasciitilde ta-wa\\
         \Rdp\textasciitilde\Fsg-go\\
    \glt `I keep going.'
  \z
\z
Sub-examples are referred to as \cref{ex:sub-a} and \cref{ex:sub-b}.

\subsection{Tables and figures}
Placeholder text before the floats.
\begin{table}
  \caption{Sample table}\label{tab:sample}
  \begin{threeparttable}
  \begin{tabular}{lll}
    \toprule
    Form & Gloss & Note \\
    \midrule
    \form{ta} & \gloss{go} & regular \\
    \form{ku} & \gloss{house} & loan\tnote{a} \\
    \bottomrule
  \end{tabular}
  \begin{tablenotes}\footnotesize
    \item[a] A table note uses a letter, not a number.
  \end{tablenotes}
  \end{threeparttable}
\end{table}

\begin{figure}
  \includegraphics[width=.6\textwidth]{example-image}
  \caption{Sample figure}\label{fig:sample}
\end{figure}

\section{Conclusion}
Placeholder.

\begin{paperappendix}
\section{Additional data}\label{app:data}
Placeholder.
\end{paperappendix}

\printglosses
\printbibliography

\end{document}
